% Vyžaduje balíčky: graphicx, float. Obrázky: analyza/analyticka_cast/obrazky/
% (\graphicspath{{analyza/analyticka_cast/obrazky/}}).
\subsection{Příklady navržených vizualizací}

Následující náhledy ukazují, jak by jednotlivé vizualizace z
tabulky~\ref{tab:vizualizace} vypadaly na reálných datech. V aplikaci budou
interaktivní, zde jde o statické ukázky vytvořené z analyzovaného souboru.

\paragraph{Mapa objektů s barvou podle atributu.} Mapa zobrazuje objekty
v obrysu a barva vyjadřuje vybraný atribut, zde způsob využití
(obrázek~\ref{fig:priklad-mapa}). Po výběru objektu se zobrazí jeho údaje.
Vedle způsobu využití lze barvou vyjádřit podlažnost nebo rok dokončení.

\begin{figure}[H]
\centering
\includegraphics[width=0.9\textwidth]{28_priklad_mapa_atribut}
\caption{Příklad mapy objektů s barvou podle způsobu využití a informací o vybraném objektu}
\label{fig:priklad-mapa}
\end{figure}

\paragraph{Hustota v síti při malém přiblížení.} Při velkém území nelze
zobrazit každý objekt zvlášť, proto se zástavba agreguje do šestiúhelníkové
sítě a barva vyjadřuje počet objektů v buňce. Příkladem je mapa na
obrázku~\ref{fig:hustota}. Při přiblížení aplikace přechází na zobrazení
jednotlivých objektů podle předchozího příkladu.

\paragraph{Časová osa a posuvník data.} Uživatel posuvníkem volí datum
a mapa ukazuje stav území k tomuto datu (obrázek~\ref{fig:priklad-osa}).
Objekty, které přibyly po začátku sledování, jsou vyznačeny odlišnou
barvou. Vybraná lokalita ukazuje novou výstavbu, která v roce 2015 ještě
neexistovala (počet objektů v území vzrostl z 317 na 410).

\begin{figure}[H]
\centering
\includegraphics[width=\textwidth]{29_priklad_casova_osa}
\caption{Příklad časové osy: stav stejného území ve třech datech}
\label{fig:priklad-osa}
\end{figure}

\paragraph{Časová řada změn, nových a zaniklých objektů.} Graf ukazuje počet
změn za měsíc. Hromadné aktualizace registru, například 01/2024 s 43\,521
změnami, jsou vyznačeny barevně a osa je logaritmická, aby bylo vidět
i běžné měsíce (obrázek~\ref{fig:priklad-rada}). Stejně lze zobrazit nové
a zaniklé objekty.

\begin{figure}[H]
\centering
\includegraphics[width=\textwidth]{30_priklad_casova_rada}
\caption{Příklad časové řady změn s odlišenými hromadnými aktualizacemi}
\label{fig:priklad-rada}
\end{figure}

\paragraph{Detail objektu s historií verzí.} Po výběru objektu se zobrazí
jeho verze v čase a změněné údaje jsou zvýrazněny (obrázek~\ref{fig:priklad-detail}).
U příkladového objektu se v roce 2024 změnilo využití z rodinného domu na bytový
dům a současně počet podlaží, bytů i zastavěná plocha.

\begin{figure}[H]
\centering
\includegraphics[width=0.95\textwidth]{31_priklad_detail_objektu}
\caption{Příklad detailu objektu s historií čtyř verzí}
\label{fig:priklad-detail}
\end{figure}

\paragraph{Srovnání obcí.} Tabulka s barevnou škálou umožňuje srovnat
vybrané obce podle několika ukazatelů najednou (obrázek~\ref{fig:priklad-srovnani}):
podílu rodinných a bytových domů, průměrné podlažnosti, stáří zástavby,
vybavení kanalizací a intenzity nové výstavby.

\begin{figure}[H]
\centering
\includegraphics[width=\textwidth]{32_priklad_srovnani_obci}
\caption{Příklad srovnání osmi největších obcí podle vybraných ukazatelů}
\label{fig:priklad-srovnani}
\end{figure}

\paragraph{Rozdělení hodnot s filtrem.} Uživatel zvolí kritéria a graf
ukazuje rozdělení vybraného atributu spolu s počtem vyhovujících objektů
(obrázek~\ref{fig:priklad-filtr}). Na příkladu jsou vybrány rodinné domy
v Brně se zastavěnou plochou 80 až 150\,m$^2$; vyhovuje 18\,269 z 30\,819
objektů.

\begin{figure}[H]
\centering
\includegraphics[width=\textwidth]{33_priklad_filtr}
\caption{Příklad filtrování s rozdělením zastavěné plochy}
\label{fig:priklad-filtr}
\end{figure}

\paragraph{Indikátor úplnosti.} Každý graf je doplněn informací o tom, u
kolika objektů je zobrazený údaj vyplněn (obrázek~\ref{fig:priklad-uplnost}).
Ukazatele se mezi obcemi liší, například podlažnost je vyplněna u 59\,\%
objektů v Kuřimi, ale u 78\,\% ve Střelicích.

\begin{figure}[H]
\centering
\includegraphics[width=\textwidth]{34_priklad_uplnost}
\caption{Příklad indikátoru úplnosti údajů podle obcí}
\label{fig:priklad-uplnost}
\end{figure}
